% ============================================================================
% WORKSHEET - Section 1.2: Recognizing Proportional Relationships
% CCSS 7.RP.A.2a
% ============================================================================
\section{Recognizing Proportional Relationships}
\setcounter{wsNumber}{0}

\begin{worksheetSkills}
\begin{itemize}
    \item Decide whether a relationship is proportional from tables, graphs, and contexts
    \item Test ratios to see whether the quotient $\frac{y}{x}$ stays constant
    \item Explain why a straight line must also pass through $(0, 0)$ to be proportional
\end{itemize}
\end{worksheetSkills}

\worksheetHeader{Proportional or Not?}

\begin{circleIt}{Circle the Proportional Relationships}

Circle every table or situation that shows a proportional relationship.

\bigskip
\textbf{A.}
\begin{center}
\begin{tabular}{|C{2cm}|C{2cm}|}
\hline
$x$ & $y$ \\
\hline
$2$ & $8$ \\
$5$ & $20$ \\
$7$ & $28$ \\
\hline
\end{tabular}
\end{center}

\textbf{B.} A taxi costs $\$4$ to start the ride and $\$2$ for each mile.

\textbf{C.}
\begin{center}
\begin{tabular}{|C{2cm}|C{2cm}|}
\hline
$x$ & $y$ \\
\hline
$3$ & $12$ \\
$6$ & $24$ \\
$9$ & $37$ \\
\hline
\end{tabular}
\end{center}

\textbf{D.} A printer makes $18$ pages every minute with no warm-up delay.

\textbf{E.} A savings account begins with $\$25$ and then grows by $\$5$ each week.
\end{circleIt}
\worksheetAnswer{The proportional relationships are A and D. In A, the ratio is constant because $8 \div 2 = 4$, $20 \div 5 = 4$, and $28 \div 7 = 4$. In D, the phrase ``$18$ pages every minute with no warm-up delay'' means the total pages can be written as $y = 18x$, which starts at zero and has a constant rate. B is not proportional because the $\$4$ start fee means the relationship does not begin at $(0, 0)$. C is not proportional because the last ratio is wrong: $37 \div 9 \neq 4$. E is not proportional because the initial $\$25$ balance is a starting amount.}

\worksheetHeader{Table Test Match}

\begin{matchIt}{Match Each Table to the Best Description}

\matchRow{Table with ratios $3, 3, 3$}{Proportional because every pair has the same unit rate}
\matchRow{Table with ratios $5, 5, 4.8$}{Not proportional because one row breaks the pattern}
\matchRow{Straight line through $(0, 0)$}{Graph shows a proportional relationship}
\matchRow{Straight line through $(0, 5)$}{Not proportional because there is a starting amount}
\matchRow{Situation with ``$\$7$ sign-up fee plus $\$3$ per hour''}{Not proportional because the graph would not pass through the origin}
\end{matchIt}
\worksheetAnswer{The matches are: ratios $3, 3, 3$ goes with ``Proportional because every pair has the same unit rate'' since the quotient never changes. Ratios $5, 5, 4.8$ goes with ``Not proportional because one row breaks the pattern'' because one mismatch is enough to show the table is not proportional. Straight line through $(0, 0)$ goes with ``Graph shows a proportional relationship'' because proportional graphs are lines through the origin. Straight line through $(0, 5)$ goes with ``Not proportional because there is a starting amount'' since the $y$-intercept is not zero. The sign-up fee situation goes with ``Not proportional because the graph would not pass through the origin'' for the same reason.}

\worksheetHeader{Find the Broken Row}

\begin{errorDetective}{Which Student Made a Wrong Claim?}

Aria says this table is proportional because the values grow by the same amount.

\errorWork{
\begin{tabular}{|c|c|}
\hline
$x$ & $y$ \\
\hline
$2$ & $7$ \\
$4$ & $12$ \\
$6$ & $17$ \\
\hline
\end{tabular}}

Is Aria correct? \answerBlank[2cm] \quad Why? \answerBlank[7cm]

\bigskip

Ben says a graph through $(1, 6)$, $(2, 12)$, and $(3, 18)$ is not proportional because it does not show $(0, 0)$ in the table.
\errorWork{``If $(0, 0)$ is missing from the table, the relationship cannot be proportional.''}

Is Ben correct? \answerBlank[2cm] \quad Why? \answerBlank[7cm]

\bigskip

Kai says a music plan with cost rule $c = 4n + 10$ is proportional because the graph is a line.
\errorWork{$c = 4n + 10$ is a straight-line equation, so it must be proportional.}

Is Kai correct? \answerBlank[2cm] \quad Why? \answerBlank[7cm]
\end{errorDetective}
\worksheetAnswer{All three students are incorrect. (1)~Aria is wrong because proportional relationships need a constant ratio, not just a constant difference. Here $7 \div 2 = 3.5$, $12 \div 4 = 3$, and $17 \div 6 \approx 2.83$, so the ratios are not equal. (2)~Ben is wrong because a table can still represent a proportional relationship even if the row for zero is not written. Since $6 \div 1 = 6$, $12 \div 2 = 6$, and $18 \div 3 = 6$, the rule is $y = 6x$, which would include $(0, 0)$. (3)~Kai is wrong because not every line is proportional. A proportional equation must have the form $y = kx$. The rule $c = 4n + 10$ has a starting fee of $10$, so its graph crosses the $y$-axis at $10$, not at zero.}

\worksheetHeader{Phone Plans and Paychecks}

\begin{mathScenario}{Decide Which Situations Are Proportional}

Review each situation.

\bigskip
\scenarioItem{Plan A}{Earn $\$14$ per hour babysitting}
\scenarioItem{Plan B}{Earn $\$14$ per hour plus a $\$10$ bonus}
\scenarioItem{Plan C}{A streaming service costs $\$6$ every month}
\scenarioItem{Plan D}{A bike rental costs $\$3$ for each half-hour with no fee to start}

\bigskip

Which plans are proportional? \answerBlank[6cm]

\bigskip

Write one sentence explaining why Plan B is not proportional. \answerBlank[8cm]

\bigskip

If Plan D lasts $2$ hours, what is the total cost, and why does that fit a proportional relationship? \answerBlank[8cm]
\end{mathScenario}
\worksheetAnswer{(1)~Plan A and Plan D are proportional. Plan A can be written $p = 14h$, so it has a constant unit rate and starts at zero. Plan D charges $\$3$ for each half-hour, so the cost is $3$ times the number of half-hours and also starts at zero. Plan B is not proportional because of the extra $\$10$ bonus, and Plan C is not proportional because it is a flat monthly fee rather than a constant ratio between two changing quantities. (2)~Plan B is not proportional because even when the hours are zero, the pay would still be $\$10$, so the graph would not pass through $(0, 0)$. (3)~Two hours is $4$ half-hours, so the cost is $4 \times 3 = \$12$. That fits a proportional relationship because the total cost is always the same rate, $\$3$ per half-hour, multiplied by the number of half-hours.}

\worksheetHeader{Explain the Origin Test}

\begin{explainIt}{Why Does the Origin Matter?}

Riley says, ``If the points form a perfect straight line, I do not need to check the origin.''
Write a reply that explains why the origin test matters for proportional relationships.

\writeLines{6}

\bigskip

Two students examine the table $(2, 10)$, $(4, 20)$, $(6, 30)$. One says it is proportional because the differences are constant. The other says it is proportional because the ratios are constant. Which explanation is stronger, and why?

\writeLines{6}
\end{explainIt}
\worksheetAnswer{(1)~A straight line is not enough because many nonproportional relationships are also straight lines. The key idea is that a proportional relationship has no starting amount, so when $x = 0$, $y$ must also equal $0$. That means the graph must pass through the origin. A line through $(0, 5)$ is straight but not proportional because the ratio $\frac{y}{x}$ changes. (2)~The stronger explanation is that the ratios are constant. In the table, $10 \div 2 = 5$, $20 \div 4 = 5$, and $30 \div 6 = 5$, so the relationship follows $y = 5x$. Constant differences can happen in many linear relationships, but constant ratios are the test for proportionality.}

\worksheetDone